% Tradução brasileira revista a partir do workpass espanhol congelado;
% a fonte permanece em source_es.
% SGA 3, Exposé VIII, tradução brasileira.
% Continuação da demonstração do Corolário 4.7 e final de 4.6.
% Autoridade: Exp8 p. local 14; leitor completo p. 630.

Isto mostra que obtemos um isomorfismo
\[
  A_0[t,t^{-1}]\longrightarrow A
\]
de álgebras \(A_0\)-graduadas que envia \(t\) em \(f\), o que conclui a
demonstração de \SGARefLink{sga3:viii:corollary:4:7}{4.7}.

No caso de
\SGARefLink{sga3:viii:proposition:4:6}{4.6},
\SGARefLink{sga3:viii:corollary:4:7}{4.7} já mostra que os
\(\mathscr A_m\), para \(m\in M\), são invertíveis. Para demonstrar a
condição \SGARefLink{sga3:viii:proposition:4:1:item:b:01}{4.1(b)}, podemos
supor, portanto, que \(\mathscr A_m\) e \(\mathscr A_{m'}\) tenham bases
\(f_m\) e \(f_{m'}\), com inversos
\[
  f_m^{-1}\in\Gamma(\mathscr A_{-m})
  \qquad\text{y}\qquad
  f_{m'}^{-1}\in\Gamma(\mathscr A_{-m'}).
\]
A multiplicação por
\[
  f_m^{-1}f_{m'}^{-1}
  \in\Gamma(\mathscr A_{-m-m'})
\]
define então um homomorfismo
\[
  \mathscr A_{m+m'}\longrightarrow
  \mathscr A_0\simeq\mathcal O_S
\]
que envia a imagem de \(f_m\otimes f_{m'}\) para a seção unidade de
\(\mathcal O_S\). No diagrama
\SGARefTarget{sga3:viii:proposition:4:6:diagram:01}
\[
\begin{tikzcd}[column sep=large]
  \mathscr A_m\otimes\mathscr A_{m'}
    \arrow[r,"u"]
    \arrow[rr,bend left=32,"w"]
  & \mathscr A_{m+m'}
    \arrow[r,"v"]
  & \mathscr A_0\simeq\mathcal O_S ,
\end{tikzcd}
\]
os morfismos \(w\) e \(v\) são, por conseguinte, epimorfismos de feixes
invertíveis e, portanto, isomorfismos; consequentemente, \(u\) é um
isomorfismo. \hfill\(\square\)
